\section{Введение}
\label{sec:Chapter0}
\index{Chapter0}

\subsection{Введение в 3D-графику и её задачи}
Современный мир невозможно представить без 3D-графики, которая прочно вошла во все сферы нашей жизни и стала одним из ключевых инструментов визуализации, моделирования и взаимодействия с данными. От интерактивных видеоигр и кинематографических спецэффектов до сложных научных симуляций, медицинских визуализаций, систем автоматизированного проектирования (САПР) и стремительно развивающихся технологий виртуальной и дополненной реальности (VR/AR),--- трехмерные изображения играют центральную роль в представлении сложных объектов, процессов и сред.

С развитием технологий и ростом вычислительных мощностей графических процессоров (GPU) постоянно возрастают требования к качеству, детализации и интерактивности 3D-изображений. Пользователи ожидают фотореалистичной графики, мгновенного отклика и возможности взаимодействия со сложными трехмерными сценами в реальном времени. Это создает постоянный вызов для разработчиков алгоритмов и систем 3D-графики, требуя поиска новых, более эффективных и гибких подходов к моделированию и, что особенно важно, к рендерингу. Особую актуальность приобретают методы, способные обеспечить высокую производительность при работе с высокодетализированными и динамически изменяющимися сценами, стремясь при этом к максимальному реализму.

\subsection{Классические подходы к представлению 3D-моделей}

На протяжении многих десятилетий доминирующим и наиболее распространенным способом представления трехмерных объектов в компьютерной графике остается использование полигональных моделей (также известных как сетки, или mesh-модели). В этом подходе поверхность объекта аппроксимируется набором плоских геометрических примитивов --- обычно треугольников или четырехугольников (квадов). Каждая вершина имеет координаты в трехмерном пространстве, а также может содержать дополнительную информацию, такую как нормали к поверхности аппроксимируемого объекта, координаты текстур и цвет.

Преимущества полигональных моделей:

\begin{itemize}
    \item
    \textbf{Простота математического описания:} Треугольник --- это простейший плоский многоугольник, математические операции с которым (интерполяция, пересечение луча и т.д.) относительно просты и хорошо изучены.
    \item
    \textbf{Хорошая поддержка аппаратным обеспечением:} Современные графические процессоры (GPU) изначально спроектированы и оптимизированы для эффективной обработки и растеризации полигональных моделей. Это позволяет достигать высокой производительности рендеринга.
    \item
    \textbf{Развитые инструменты и стандарты:} Существует огромное количество программного обеспечения для 3D-моделирования (например, Blender, Autodesk Maya, 3ds Max) и межплатформенных форматов (например, OBJ, FBX, GLTF), которые активно используются и постоянно развиваются индустрией.
    \item
    \textbf{Простота визуализации:} Прямая растеризация полигонов через стандартный графический конвейер (pipeline) обеспечивает высокую скорость вывода изображения на экран.
\end{itemize}

Однако, несмотря на свои преимущества, полигональные модели имеют ряд существенных ограничений, особенно при работе со сложными, органическими или динамически изменяющимися поверхностями:

\begin{itemize}
    \item
    \textbf{Топологические ограничения:} Слишком плотная или, наоборот, разреженная полигональная сетка может создавать проблемы при моделировании и деформации. Создание плавных, округлых форм требует большого количества полигонов, что увеличивает объем данных.
    \item
    \textbf{Проблема уровня детализации (Level of Detail, LOD):} Для эффективного рендеринга объектов на разных расстояниях от камеры (или с разной степенью важности) необходимо иметь несколько версий одной и той же модели с различным числом полигонов. Управление этими версиями и переключение между ними требует дополнительных усилий и усложняет конвейер разработки.
    \item
    \textbf{Сложность описания органических форм и динамических изменений:} Точное представление объектов со сложной, изменяющейся топологией (например, жидкости, облака, мягкие тела, волосы) с помощью фиксированной полигональной сетки является крайне сложной или невозможной задачей без постоянного перестроения сетки.
    \item
    \textbf{Булевы операции:} Выполнение таких операций, как объединение, пересечение или вычитание полигональных моделей, является нетривиальной задачей и часто приводит к созданию некачественной (non-manifold, с самопересечениями) или избыточной геометрии, требующей ручной доработки.
    \item
    \textbf{Ограничения по памяти:} Для достижения высокой детализации требуется экспоненциально увеличивать количество полигонов, что приводит к большим объемам данных и потенциальным проблемам с хранением, передачей и обработкой.
\end{itemize}

Эти ограничения обусловили поиск альтернативных способов представления 3D-моделей, которые могли бы обеспечить большую гибкость и эффективность, особенно в условиях возрастающих требований к сложности и динамичности трехмерных сцен.

\subsection{Неявные поверхности: функции расстояния}
Ограничения полигонального представления, особенно в контексте моделирования сложных органических форм, выполнения булевых операций и потребности в адаптивном уровне детализации, стимулировали развитие альтернативных методов представления трехмерных объектов. Одним из таких направлений стало использование неявных поверхностей (Implicit Surfaces).

В отличие от явного (explicit) представления, где поверхность описывается параметрически в виде набора точек, линий или полигонов (например, полигональные сетки, NURBS-кривые \cite{piegl2012nurbs}), неявная поверхность определяется как множество точек (x,y,z), для которых некоторая функция F(x,y,z) обращается в ноль (то есть F(x,y,z)=0).

Особое значение среди неявных поверхностей получили функции расстояния (Distance Functions), а в частности, знаковые функции расстояния (Signed Distance Fields, SDFs) \cite{Osher2003}. SDF представляет собой функцию D(x,y,z), которая для любой точки в пространстве возвращает кратчайшее евклидово расстояние до ближайшей точки поверхности объекта. При этом знак функции указывает, находится ли точка внутри объекта (отрицательное значение) или снаружи (положительное значение). Нулевой уровень этой функции (D(x,y,z)=0) точно определяет поверхность объекта.

Ключевые преимущества SDFs:

\begin{itemize}
    \item
    \textbf{Компактное представление:} Сложные топологии и детализированные поверхности могут быть описаны относительно простой аналитической функцией или дискретным набором значений, хранимых в иерархических структурах данных (например, окдеревьях, воксельных сетках), вместо огромного числа полигонов. Это обеспечивает высокую экономию памяти для хранения.
    \item
    \textbf{Гибкость в создании и модификации:} SDFs позволяют легко выполнять булевы операции (объединение, пересечение, вычитание) над сложными формами простыми математическими операциями. Это значительно упрощает процедурное моделирование и создание сложных объектов.
    \item
    \textbf{``Бесконечная'' детализация:} Поскольку поверхность определяется функцией, а не дискретным набором данных, возможно получение любого уровня детализации ``на лету'' без потери информации. Для увеличения детализации достаточно просто увеличить количество ``запросов'' к функции расстояния.
    \item
    \textbf{Простота получения нормалей:} В любой точке поверхности вектор нормали может быть легко вычислен как градиент функции расстояния.
    \item
    \textbf{Наличие информации о расстоянии до поверхности:} Знать расстояние до поверхности полезно для определения столкновений, расчета толщины объекта, сглаживания краев и других применений.
    \item
    \textbf{Использование в различных областях:} SDFs находят широкое применение в процедурной генерации контента (пейзажи, персонажи, текстуры), симуляции жидкостей и мягких тел, 3D-печати (для проверки корректности геометрии), быстрого прототипирования, а также для хранения геометрии в задачах машинного обучения.
\end{itemize}

Несмотря на свои значительные преимущества, использование SDFs сопряжено с рядом фундаментальных проблем, особенно в контексте интерактивной визуализации:

\begin{itemize}
    \item
    \textbf{Несовместимость с традиционным полигональным графическим конвейером:} Стандартные GPU-конвейеры разработаны для рендеринга полигональных моделей. Прямой рендеринг SDFs требует специализированных алгоритмов, таких как трассировка лучей, которые могут быть вычислительно затратными.
    \item
    \textbf{Сложность эффективной конвертации в полигональную сетку в реальном времени:} Хотя методы вроде Marching Cubes позволяют преобразовать SDF в полигоны, они часто требуют дискретизации пространства (вокселизации) и сопряжены с проблемами потери детализации или генерации избыточного количества геометрии. Создание адаптивной и высококачественной полигональной сетки из SDF в реальном времени остается сложной задачей.
\end{itemize}

Эти вызовы подчеркивают необходимость разработки новых, высокопроизводительных и адаптивных методов рендеринга SDFs, способных эффективно использовать преимущества данного представления и соответствовать требованиям современных графических приложений.

\subsection{Актуальность диссертации}
В последние годы, с бурным развитием графических процессоров и появлением новых аппаратных архитектур, таких как Mesh Shaders \cite{NVIDIA_MeshShaders} в современных GPU, открываются принципиально новые возможности для переосмысления процессов рендеринга неявных поверхностей. Mesh Shaders позволяют программировать конвейер генерации геометрии непосредственно на GPU, обеспечивая беспрецедентную гибкость и управляемость на ранних стадиях графического конвейера. Это создает уникальную возможность для реализации высокоэффективных алгоритмов динамической триангуляции неявных поверхностей, адаптированных под современное аппаратное обеспечение.

Таким образом, актуальность данной магистерской диссертации определяется острой потребностью в разработке и оптимизации методов рендеринга 3D-моделей, заданных функциями расстояния, которые могли бы:

\begin{itemize}
    \item
     Эффективно использовать современные аппаратные возможности GPU, в частности, Mesh Shaders, для максимально производительной и адаптивной генерации геометрии ``на лету''.
    \item
     Снизить вычислительные затраты и объем генерируемой геометрии за счет адаптивной триангуляции видимых участков поверхности.
\end{itemize}

Разработка такого метода, базирующегося на динамической триангуляции и специализированных GPU-конвейерах, позволит преодолеть ограничения существующих подходов и сделает интерактивную визуализацию сложных SDF-моделей более доступной и производительной, открывая новые горизонты для их применения в областях, где сейчас доминирует полигональное представление.